\chapter*{Abstract}
\addstarredchapter{Abstract}

\initial{C}onversion from mild cognitive impairment (MCI) to dementia is a central clinical and public-health question. Multimodal machine-learning models often report AUC above 0.8, yet about two thirds are trained on ADNI and only about 10\% are externally validated; longitudinal leakage and unstable phenotypes are common. This two-month internship (24 July -- 24 September 2026) at UMMISCO, in collaboration with the Glenn Biggs Institute, targets a leakage-resistant 2- and 5-year conversion-risk model, external validation, and a still-useful \emph{minimal data package}.

The contribution is not another internal AUC. It is: (i) a map of NACC, ADNI and Framingham access routes; (ii) a literature review (CARD/NIH, MOIRA, Kolachalama, leakage, PROBAST); (iii) an analysis plan in the Fongang-lab format (three aims, locked phenotypes, nested models 1--4, calibration, DCA, SHAP, fairness). Working assumption: mentor-held extracts under lab DUAs. The analysis unit is the participant at an MCI index visit; all of that person's visits stay in one split. Incomplete profiles (no PET or CSF) are kept.

\paragraph{Keywords}: MCI, dementia conversion, multimodal learning, data leakage, external validation, NACC, ADNI, Framingham, minimal data package.
